\chapter{RX: the Machine's REXX}\label{cha:rx}

\pth{/STARTUP.BAT} is a list of commands, and a list of commands is
enough right up to the moment something has to be decided. RX is the
answer to that: a REXX for this machine, in the Amiga's tradition rather
than the mainframe's.

\begin{type}
RX HELLO
\end{type}

\noindent runs \texttt{HELLO.RX}. So does typing \verb|HELLO| on its own,
because the shell's rule for an unknown word --- try it as a program on
disk --- takes a script as readily as a \texttt{.prg}. The scripts that
come with the machine are in \pth{/LANG/RX}.

\section{Why REXX, and what ARexx was}
ARexx on the Amiga was two things at once. It was the REXX language ---
readable, typeless, and forgiving in the way a scripting language for
people who are not programmers has to be. And it was the \emph{port}
system: every serious Amiga program published a port, so ten lines of
REXX could drive a paint program and a terminal in the same script. The
second half is why anybody remembers it.

Both halves are here. The language is \pth{demo/rexx.c}, 37\,KB loaded
at \texttt{\$0800} because it wants nearly the whole program area, and it
is a real one: strings and whole numbers, \verb|PARSE| with its
templates, \verb|DO| in all its forms, \verb|SELECT|, internal functions
with \verb|PROCEDURE EXPOSE| and recursion, compound (stem) variables,
\verb|SIGNAL|, and about fifty built-in functions.

The thing that makes it glue is REXX's oldest idea: \textbf{a clause the
interpreter does not recognise is not an error, it is a command}, sent
to whatever the current environment is.

\begin{type}
ADDRESS COMMAND
'DIR APPS'
IF RC = 0 THEN SAY 'that worked'
\end{type}

\section{The three environments}
\begin{center}
\small
\begin{tabular}{@{}l >{\raggedright\arraybackslash}p{0.62\linewidth}@{}}
\toprule
\textbf{ADDRESS} & \textbf{Sends the line to} \\
\midrule
\texttt{COMMAND} & K/OS's own shell --- exactly as if you had typed it. The result comes back in \verb|RC| \\
\texttt{TUBE} & the co-processor's shell, with its output brought back in \verb|RESULT| \\
\emph{name} & a program's port (below) \\
\bottomrule
\end{tabular}
\end{center}

\section{Ports: the file is the mailbox}
The Amiga's ports were message ports between tasks that were all running
at once. This machine runs one program at a time, so that arrangement
could not be copied --- and the substitute is the obvious one: a file.

A script writes its request to \pth{/LANG/RX/MAIL.CMD}, \verb|SWAP|s to
the program with that file as its one argument, and reads the reply from
\pth{/LANG/RX/MAIL.RPL}: a result code on the first line and the text
after it. A program joins the scheme by honouring one argument, which is
about twenty lines of work.

\verb|CHESS| is the worked example, and it plays a whole game this way:
\verb|NEW|, \verb|MOVE|, \verb|GO|, \verb|LEVEL|, \verb|BOARD|,
\verb|FEN| and \verb|STATUS|. Its position lives in
\pth{/APPS/CHESS/PORT.GAM}, because a program swapped in starts fresh
every time and the file is its memory. \texttt{CHESS.RX} in
\pth{/LANG/RX} plays against the engine and draws the board itself from
what comes back.

\section{What is in \texttt{/LANG/RX}}
A language with nothing written in it proves nothing, so:

\begin{center}
\small
\begin{tabular}{@{}l >{\raggedright\arraybackslash}p{0.58\linewidth}@{}}
\toprule
\texttt{HELLO.RX} & the smallest one \\
\texttt{MACHINE.RX} & the machine reading its own registers \\
\texttt{HOSTINFO.RX} & the host asked about itself, over the Tube \\
\texttt{WC.RX} & a word count \\
\texttt{SUMFILE.RX} & a column adder \\
\texttt{GUESS.RX} & a guessing game \\
\texttt{WUMPUS.RX} & Hunt the Wumpus, in 150 lines \\
\texttt{CHESS.RX} & a game of chess played through the port \\
\bottomrule
\end{tabular}
\end{center}

\begin{aside}
\textbf{Games are good tests.} Wumpus alone found three real faults in
the interpreter on the day it was written: \verb|EXPOSE| of a stem was
\emph{assigning} it, which silently threw away every member of the
caller's map; \verb|LEAVE| stopped at the first \verb|END| it met rather
than the loop's, so leaving from inside a \verb|SELECT| went nowhere;
and a plain \verb|DO| group was counting as a loop. None of them would
have shown up in a script that only ran commands in order.
\end{aside}

\begin{aside}
\textbf{One thing to know about running programs from a script.} A
program started by the shell loads at \texttt{\$6000} --- through the
middle of the interpreter. So RX asks the filesystem whether the first
word of a command names a file: a built-in shell command runs directly
and its output stays on the screen, and a program is swapped in and out
around the interpreter. That is also why RX uses \verb|SWAP -k|
(Chapter~\ref{cha:shell}): the ordinary \verb|SWAP| would restore the
screen and throw away what the program had just printed for the script
to read.
\end{aside}
